%============================================================
% HALAMAN PENGESAHAN
% Halaman ini selalu dicetak tanpa watermark.
%============================================================
\cleardoublepage
\pagestyle{empty} % Memastikan seluruh halaman pengesahan tanpa nomor/header

\DisableDocumentWatermark

\begin{center}
  {\fontsize{20}{0}\bfseries HALAMAN PENGESAHAN\par}
  \vspace{0.8cm}

  \begin{minipage}{1\textwidth}
    \centering
    \fontsize{18}{16}\bfseries\selectfont\MakeUppercase{\JudulPA}
  \end{minipage}
  \par

  \vspace{0.6cm}

  {\fontsize{14}{17}\selectfont Oleh:\par}
  \vspace{0.1cm}

  {\fontsize{14}{17}\bfseries
    \NamaMahasiswa\par
  NRP. \NRP\par}

  \vspace{0.6cm}

  \begin{spacing}{1}
    {\fontsize{14}{17}\bfseries
      Proyek Akhir ini digunakan sebagai salah satu syarat untuk\\
      memperoleh Gelar \GelarLulus\\
      di \NamaInstitusi\\
    \TahunLulus\par}
  \end{spacing}
\end{center}

\vspace{0.2cm}
\begin{center}
  {\fontsize{14}{17}\bfseries Disetujui oleh:}
\end{center}
\vspace{0.2cm}

% Daftar pembimbing dan penguji beserta kolom tanda tangan
\noindent
{\fontsize{14}{17}\selectfont
  \begin{tabular}{@{}l c l c@{}}
    Pembimbing 1 & : & \NamaPembimbingSatu   & (\hspace{2.5cm}) \\
    &   & NIP. \NIPPembimbingSatu & \\[3pt]
    \ifcsvoid{NamaPembimbingDua}{}{%
      Pembimbing 2 & : & \NamaPembimbingDua & (\hspace{2.5cm}) \\
      & & \ifcsvoid{NIPPembimbingDua}{}{NIP. \NIPPembimbingDua} & \\[3pt]
    }

    % Baris kondisional Pembimbing 3 (Akademik)
    \ifcsvoid{NamaPembimbingTiga}{}{%
      Pembimbing 3 & : & \NamaPembimbingTiga   & (\hspace{2.5cm}) \\
      & & \ifcsvoid{NIPPembimbingTiga}{}{NIP. \NIPPembimbingTiga} & \\[3pt]
    }

    % Baris kondisional Pembimbing Industri
    \ifcsvoid{NamaPembimbingIndustri}{}{%
      Pembimbing Industri & : & \NamaPembimbingIndustri & (\hspace{2.5cm}) \\
      & & \IdentitasPembimbingIndustri & \\[3pt]
    }

    Penguji 1    & : & \NamaPengujiSatu     & (\hspace{2.5cm}) \\
    &   & NIP. \NIPPengujiSatu   & \\[3pt]
    Penguji 2    & : & \NamaPengujiDua      & (\hspace{2.5cm}) \\
    &   & NIP. \NIPPengujiDua    & \\
  \end{tabular}
}

\vspace{0.3cm}

% 2. Blok Mengetahui dibungkus dalam minipage
%    Jika ruang tidak cukup, SELURUH kotak ini akan pindah ke halaman 2.
\noindent
\begin{minipage}{\textwidth}
  \begin{center}
    {\fontsize{14}{17}\bfseries Mengetahui,}\par
    {\fontsize{14}{17}\selectfont
      Ketua Program Studi Sarjana Terapan \NamaProdi\par
    \NamaInstitusi}\par
    \vspace{1.5cm}
    {\fontsize{14}{17}\selectfont
      \underline{\NamaKetuaProdi}\par
    NIP. \NIPKetuaProdi\par}
  \end{center}
\end{minipage}

% Paksa render halaman terakhir sebelum masuk ke bagian berikutnya.
\clearpage
